% GENERATED FILE. Edit canonical lessons, then run npm run generate:books.
% canonical-source-hash: fnv1a64:13778bfd61351a59
% canonical-lessons: MW-C254-vaid, MW-C254-goli, MW-C254-sujno, MW-C254-aram, MW-C254-alas

\chapter{At the Healer's}
\label{ch:mw-at-the-healer-s}
\hlchaptermodality{254}

\begin{chapteropening}
\textbf{By the end of this chapter:} \emph{I can say what the healer gave me and what I need.}

Complete the last lesson of chapter 254: I can say what the healer gave me and what I need.
\end{chapteropening}

\section[\texorpdfstring{\mw{वैद} (vaid)}{वैद (vaid)}]{\mw{वैद} (vaid) — a traditional healer; a physician}

\label{lesson:MW-C254-vaid}



\begin{quote}
Before the new one: say the Marwadi for to throb, then the Marwadi for to ache.
\end{quote}

\subsection*{You'll want to know: \mw{वैद}}
\textbf{\mw{वैद}} — \emph{vaid} — \textquotedblleft{}a traditional healer; a physician\textquotedblright{}.

\begin{quote}
\mw{गांव} \mw{मांय} \mw{एक} \mw{वैद} \mw{है।}
\end{quote}

\begin{quote}
\emph{gā̃v mā̃y ek vaid hai.}
\end{quote}

\begin{quote}
There is a healer in the village.
\end{quote}

A masculine noun.

\begin{cousinweb}
\textbf{\mw{वैद}} (\emph{vaid}) goes back to Sanskrit \emph{vaidya}, \textquotedblleft{}learned; a physician\textquotedblright{} (Monier-Williams), with its last syllable worn away in speech.
\end{cousinweb}

\begin{grammarlens}[title={What you've built}]
One more word for this chapter.
\end{grammarlens}

\subsection*{Your turn}
\begin{itemize}
\raggedright
  \item \emph{Say it:} \emph{vaid}
  \item \emph{Say it:} \emph{vaid}, once more
  \item \emph{Say it:} \emph{gā̃v mā̃y ek vaid hai.} — There is a healer in the village.
  \item \emph{Recall:} say the Marwadi for to throb, then the Marwadi for to ache, then say \emph{vaid} again
\end{itemize}

\begin{tcolorbox}[breakable,colback=teal!4,colframe=teal!35!black,title={Before you move on}]
What does \mw{वैद} mean? (\textquotedblleft{}A traditional healer; a physician\textquotedblright{}.) Say it once more.
\end{tcolorbox}
\section[\texorpdfstring{\mw{गोळी} (goḷī)}{गोळी (goḷī)}]{\mw{गोळी} (goḷī) — a pill; a tablet}

\label{lesson:MW-C254-goli}



\begin{quote}
Before the new one: say the Marwadi for to ache, then the Marwadi for a traditional healer.
\end{quote}

\subsection*{You'll want to know: \mw{गोळी}}
\textbf{\mw{गोळी}} — \emph{goḷī} — \textquotedblleft{}a pill; a tablet\textquotedblright{}.

\begin{quote}
\mw{म्हैं} \mw{दिनूगै} \mw{गोळी} \mw{खाई।}
\end{quote}

\begin{quote}
\emph{mhaĩ dinūgai goḷī khāī.}
\end{quote}

\begin{quote}
I took a pill early in the morning.
\end{quote}

A feminine noun.

\begin{cousinweb}
\textbf{\mw{गोळी}} (\emph{goḷī}) is inherited from Sanskrit \emph{gola}, \textquotedblleft{}a ball\textquotedblright{}. CDIAL gives no Marwari form in this entry; it lists Gujarati \emph{goḷī}, \textquotedblleft{}pill\textquotedblright{} (CDIAL 4321).
\end{cousinweb}

\begin{grammarlens}[title={What you've built}]
One more word for this chapter.
\end{grammarlens}

\subsection*{Your turn}
\begin{itemize}
\raggedright
  \item \emph{Say it:} \emph{goḷī}
  \item \emph{Say it:} \emph{goḷī}, once more
  \item \emph{Say it:} \emph{mhaĩ dinūgai goḷī khāī.} — I took a pill early in the morning.
  \item \emph{Recall:} say the Marwadi for to ache, then the Marwadi for a traditional healer, then say \emph{goḷī} again
\end{itemize}

\begin{tcolorbox}[breakable,colback=teal!4,colframe=teal!35!black,title={Before you move on}]
What does \mw{गोळी} mean? (\textquotedblleft{}A pill; a tablet\textquotedblright{}.) Say it once more.
\end{tcolorbox}
\section[\texorpdfstring{\mw{सूजणो} (sūjṇo)}{सूजणो (sūjṇo)}]{\mw{सूजणो} (sūjṇo) — to swell up}

\label{lesson:MW-C254-sujno}



\begin{quote}
Before the new one: say the Marwadi for a traditional healer, then the Marwadi for a pill.
\end{quote}

\subsection*{You'll want to know: \mw{सूजणो}}
\textbf{\mw{सूजणो}} — \emph{sūjṇo} — \textquotedblleft{}to swell up\textquotedblright{}.

\begin{quote}
\mw{म्हारो} \mw{पग} \mw{सूज्यो।}
\end{quote}

\begin{quote}
\emph{mhāro pag sūjyo.}
\end{quote}

\begin{quote}
My foot swelled up.
\end{quote}

Past (it …): \mw{वो} \mw{गोडो} \mw{सूज्यो}, \emph{vo goḍo sūjyo}.

\begin{cousinweb}
\textbf{\mw{सूजणो}} (\emph{sūjṇo}) is inherited from Sanskrit \emph{śūyate}, \textquotedblleft{}is swollen\textquotedblright{}. CDIAL gives no Marwari form in this entry; it lists Hindi \emph{sūjnā} (CDIAL 12568).
\end{cousinweb}

\begin{grammarlens}[title={What you've built}]
One more word for this chapter.
\end{grammarlens}

\subsection*{Your turn}
\begin{itemize}
\raggedright
  \item \emph{Say it:} \emph{sūjṇo}
  \item \emph{Say it:} \emph{sūjṇo}, once more
  \item \emph{Say it:} \emph{mhāro pag sūjyo.} — My foot swelled up.
  \item \emph{Recall:} say the Marwadi for a traditional healer, then the Marwadi for a pill, then say \emph{sūjṇo} again
\end{itemize}

\begin{tcolorbox}[breakable,colback=teal!4,colframe=teal!35!black,title={Before you move on}]
What does \mw{सूजणो} mean? (\textquotedblleft{}To swell up\textquotedblright{}.) Say it once more.
\end{tcolorbox}
\section[\texorpdfstring{\mw{आराम} (ārām)}{आराम (ārām)}]{\mw{आराम} (ārām) — rest}

\label{lesson:MW-C254-aram}



\begin{quote}
Before the new one: say the Marwadi for a pill, then the Marwadi for to swell up.
\end{quote}

\subsection*{You'll want to know: \mw{आराम}}
\textbf{\mw{आराम}} — \emph{ārām} — \textquotedblleft{}rest\textquotedblright{}.

\begin{quote}
\mw{आज} \mw{आराम} \mw{करो।}
\end{quote}

\begin{quote}
\emph{āj ārām karo.}
\end{quote}

\begin{quote}
Rest today.
\end{quote}

A masculine noun.

\begin{cousinweb}
\textbf{\mw{आराम}} (\emph{ārām}) is usually traced to Persian \emph{ārām}, \textquotedblleft{}rest, ease\textquotedblright{}.
\end{cousinweb}

\begin{grammarlens}[title={What you've built}]
One more word for this chapter.
\end{grammarlens}

\subsection*{Your turn}
\begin{itemize}
\raggedright
  \item \emph{Say it:} \emph{ārām}
  \item \emph{Say it:} \emph{ārām}, once more
  \item \emph{Say it:} \emph{āj ārām karo.} — Rest today.
  \item \emph{Recall:} say the Marwadi for a pill, then the Marwadi for to swell up, then say \emph{ārām} again
\end{itemize}

\begin{tcolorbox}[breakable,colback=teal!4,colframe=teal!35!black,title={Before you move on}]
What does \mw{आराम} mean? (\textquotedblleft{}Rest\textquotedblright{}.) Say it once more.
\end{tcolorbox}
\section[\texorpdfstring{\mw{आळस} (āḷas)}{आळस (āḷas)}]{\mw{आळस} (āḷas) — sluggishness; laziness}

\label{lesson:MW-C254-alas}



\begin{quote}
Before the new one: say the Marwadi for to swell up, then the Marwadi for rest.
\end{quote}

\subsection*{You'll want to know: \mw{आळस}}
\textbf{\mw{आळस}} — \emph{āḷas} — \textquotedblleft{}sluggishness; laziness\textquotedblright{}.

\begin{quote}
\mw{ताव} \mw{मांय} \mw{घणो} \mw{आळस} \mw{लागै} \mw{है।}
\end{quote}

\begin{quote}
\emph{tāv mā̃y ghaṇo āḷas lāgai hai.}
\end{quote}

\begin{quote}
With a fever you feel very sluggish.
\end{quote}

A masculine noun.

\begin{cousinweb}
\textbf{\mw{आळस}} (\emph{āḷas}) is inherited from Sanskrit \emph{ālasya}, \textquotedblleft{}sloth\textquotedblright{}. CDIAL gives no Marwari form in this entry; it lists Gujarati \emph{āḷas}, with this book's \mw{ळ} (CDIAL 1371).
\end{cousinweb}

\begin{grammarlens}[title={What you've built}]
One more word for this chapter.
\end{grammarlens}

\subsection*{Your turn}
\begin{itemize}
\raggedright
  \item \emph{Say it:} \emph{āḷas}
  \item \emph{Say it:} \emph{āḷas}, once more
  \item \emph{Say it:} \emph{tāv mā̃y ghaṇo āḷas lāgai hai.} — With a fever you feel very sluggish.
  \item \emph{Recall:} say the Marwadi for to swell up, then the Marwadi for rest, then say \emph{āḷas} again
\end{itemize}

\begin{tcolorbox}[breakable,colback=teal!4,colframe=teal!35!black,title={Before you move on}]
What does \mw{आळस} mean? (\textquotedblleft{}Sluggishness; laziness\textquotedblright{}.) Say it once more.
\end{tcolorbox}
